\documentclass[11pt,a4paper]{article}
\usepackage[top=1.7cm,bottom=1.5cm,left=1.5cm,right=1.5cm,headsep=8pt,footskip=16pt]{geometry}
\usepackage{enumitem}
\usepackage{amssymb}
\usepackage{array}
\usepackage[table]{xcolor}
\usepackage{fancyhdr}
\usepackage{lastpage}
\usepackage{multicol}
\usepackage{polyglossia}
\setdefaultlanguage[calendar=gregorian, numerals=maghrib]{arabic}
\setotherlanguage{english}
\IfFontExistsTF{Amiri}{%
  \newfontfamily\arabicfont[Script=Arabic, Scale=1.08]{Amiri}
}{%
  \newfontfamily\arabicfont[Script=Arabic, Scale=1.08, Path=../fonts/, Extension=.ttf,
    UprightFont=*-Regular, BoldFont=*-Bold, ItalicFont=*-Italic,
    BoldItalicFont=*-BoldItalic]{Amiri}
}
\IfFontExistsTF{AlBattar}{%
  \newfontfamily\khatt[Script=Arabic, Scale=1.15]{AlBattar}
}{%
  \newfontfamily\khatt[Script=Arabic, Scale=1.15, Path=../fonts/,
    BoldFont=ae_AlBattar.ttf]{ae_AlBattar.ttf}
}
\renewcommand{\headrulewidth}{0.8pt}
\renewcommand{\footrulewidth}{0.8pt}
\fancyhf{}
\rhead{\khatt المدرسة العليا للأساتذة ببوزريعة}
\lhead{\khatt ملحقة بومرداس}
\lfoot{\khatt الأستاذ: د. بن عبد الرزاق يوسف}
\cfoot{\thepage/\pageref{LastPage}}
\pagestyle{fancy}
\parindent 0pt
\setlength{\columnsep}{0.6cm}
\setlength{\columnseprule}{0.4pt}
\setlist[enumerate,1]{label=\textbf{\arabic*.}, leftmargin=*, itemsep=5pt, topsep=2pt}
\begin{document}

\linespread{1.08}\selectfont
\begin{center}
{\LARGE\khatt\textbf{الاختبار العشرون في الذكاء الاصطناعي}}\\[4pt]
{\normalsize وضعيات إشكالية ــ أسئلة متعددة الاختيارات ــ الفصول من 1 إلى 6}
\end{center}
\vspace{-2pt}\rule{\textwidth}{0.6pt}\\[2pt]
\textbf{التعليمة:} يتضمن الاختبار خمس وضعيات إشكالية. اقرأ كل وضعية بعناية، ثم اختر التحليل أو التصرف الأنسب
من بين الاختيارات الخمسة؛ لكل وضعية إجابة صحيحة واحدة فقط.
اكتب \textbf{حرف} الإجابة الصحيحة في جدول الإجابات في آخر الاختبار.
لكل وضعية 4 نقاط (المجموع 20 نقطة).
\vspace{4pt}

\noindent
{\large\khatt\textbf{الوضعية 1}}\hfill{\small(الفصل 1)}\par\nopagebreak\vspace{3pt}
\setlength{\fboxsep}{7pt}\colorbox{gray!10}{\parbox{\dimexpr\linewidth-14pt}{طلب أستاذ من أداة توليدية «قصيدة عن الأم» لمسابقة إلقاء، فلاحظ زميل أن بعض أبياتها تكاد تطابق أبياتًا مشهورة لشاعر معاصر. ويرى الأستاذ أن «الأداة تولّد نصوصًا جديدة دائمًا». وطلب الأستاذ من التلميذ الفائز في الانتقاء الأولي أن يحفظ القصيدة، والمسابقة ستُنقل على الصفحة الرسمية للمؤسسة. ما التقدير الأدق؟}}\par\nopagebreak
\begin{itemize}[leftmargin=1.0cm, labelsep=6pt, itemsep=3pt, topsep=4pt, parsep=0pt]
  \item[\textbf{أ)}] زميله مخطئ بالضرورة، لأن الأدوات التوليدية لا يمكن أن تنتج أي مقطع يشبه نصًا موجودًا، فكل ما تنتجه جديد حرفيًا، والتشابه الملاحظ وهم ناتج عن شيوع الموضوع.
  \item[\textbf{ب)}] يطلب من الأداة تغيير الكلمات المتشابهة فقط، ثم يستعمل القصيدة دون نسبة، لأن تغيير بعض الكلمات يكفي لتحويل الأبيات المطابقة إلى إبداع جديد لا يحتاج إلى توثيق، وهذا ما يفعله كثيرون.
  \item[\textbf{ج)}] الأداة تنسخ دائمًا نصوصًا موجودة ولا تولّد أي شيء جديد، فكل القصيدة منقولة بالضرورة، والحل البحث عن الشاعر الأصلي لكل بيت فيها، حتى الأبيات التي لا تشبه أي نص، ثم يحذف كل الأبيات المشكوك فيها.
  \item[\textbf{د)}] الأداة تولّد نصوصًا جديدة في الغالب، لكنها قد تعيد إنتاج مقاطع قريبة جدًا من نصوص شائعة في بيانات تدريبها؛ فيتحقق من أصالة الأبيات قبل استعمالها، وينسب ما يطابق نصًا معروفًا إلى صاحبه.
  \item[\textbf{هـ)}] التشابه مقبول ما دامت الأداة هي التي كتبت، لأن النص المولَّد لا يُنسب إلى أحد، فيستعمل القصيدة في المسابقة دون تحقق أو نسبة، ولو طابقت أبياتها نصًا مشهورًا.
\end{itemize}
\par\vspace{6pt}
\noindent
{\large\khatt\textbf{الوضعية 2}}\hfill{\small(الفصل 2)}\par\nopagebreak\vspace{3pt}
\setlength{\fboxsep}{7pt}\colorbox{gray!10}{\parbox{\dimexpr\linewidth-14pt}{لاحظت أستاذة أن أداة ذكاء اصطناعي تجيب بالثقة نفسها والأسلوب الحاسم نفسه سواء أصابت أم أخطأت، حتى في أسئلة دقيقة عن تواريخ ونسبة أبيات. وسألها تلاميذها: «كيف نعرف متى تكون الأداة متأكدة؟». وتريد أن تعلّمهم قاعدة عملية يطبقونها في كل استعمال للأداة، في الدراسة وخارجها. أي جواب هو الأدق؟}}\par\nopagebreak
\begin{itemize}[leftmargin=1.0cm, labelsep=6pt, itemsep=3pt, topsep=4pt, parsep=0pt]
  \item[\textbf{أ)}] النموذج تعلّم إنتاج نصوص سلسة واثقة الصياغة، ولا تعكس نبرته بالضرورة صحة المعلومة؛ فلا تُتخذ الثقة في الأسلوب دليلًا، ويُتحقق من المعلومات الدقيقة في مصادر موثوقة مهما بدت الإجابة حاسمة.
  \item[\textbf{ب)}] نعرف ذلك من طول الإجابة، فالأداة تطيل عندما تكون متأكدة وتختصر عندما تشك، ولذلك تُعتمد الإجابات الطويلة دون تحقق، ويُتحقق من الإجابات القصيرة فقط، فهذا مؤشر بسيط يفهمه كل التلاميذ.
  \item[\textbf{ج)}] نسأل الأداة «هل أنتِ متأكدة؟»، فإن أجابت بنعم كانت الإجابة صحيحة بالضرورة، لأن النماذج تعرف بدقة حدود معرفتها، وتصرّح بها بصدق عند السؤال.
  \item[\textbf{د)}] الأداة متأكدة دائمًا لأنها لا تقدّم إلا ما تعرفه، والأسلوب الحاسم دليل على الصحة، فإذا أخطأت مرة فالخطأ في المصدر الذي تعلّمت منه لا في الأداة، ولذلك لا تحتاج إلى أي تحقق.
  \item[\textbf{هـ)}] نعرف ذلك من استعمال الأداة لعبارات مثل «بالتأكيد» و«بلا شك»، فهي لا تستعملها إلا مع المعلومات المؤكدة، فتكفي هذه العبارات للحكم على صحة الإجابة.
\end{itemize}
\par\vspace{6pt}
\noindent
{\large\khatt\textbf{الوضعية 3}}\hfill{\small(الفصل 4)}\par\nopagebreak\vspace{3pt}
\setlength{\fboxsep}{7pt}\colorbox{gray!10}{\parbox{\dimexpr\linewidth-14pt}{تبيّن لأستاذ أن عددًا من طلبته استعانوا بأدوات الذكاء الاصطناعي في اختبار إلكتروني منزلي يتكون من أسئلة تعريفات واسترجاع معلومات، فتطابقت إجاباتهم تقريبًا. ويريد معالجة الأمر حاضرًا ومستقبلًا بإنصاف. واحتجّ بعض الطلبة بأن الأستاذ لم يحدد ما المسموح به، وطالب آخرون بمعاقبة الجميع دون استثناء، والإدارة تنتظر تقريره. ما المعالجة الأنسب؟}}\par\nopagebreak
\begin{itemize}[leftmargin=1.0cm, labelsep=6pt, itemsep=3pt, topsep=4pt, parsep=0pt]
  \item[\textbf{أ)}] يتجاهل الأمر تمامًا، لأن استعمال الأدوات في الاختبارات المنزلية أمر واقع، ويحتفظ بأسئلة الاسترجاع نفسها، لأنها الأسهل في التصحيح، دون أي قواعد جديدة.
  \item[\textbf{ب)}] يستعمل برنامج كشف النصوص المولَّدة ويعاقب كل من تتجاوز نسبته 50\%، دون محاورة الطلبة، لأن الكاشف دقيق، ويحتفظ بأسئلة الاسترجاع نفسها في الاختبارات القادمة، فالكاشف يوفر عليه وقت التحقق ويحمي سمعة الاختبار.
  \item[\textbf{ج)}] يمنح كل الطلبة صفرًا دون تحقق أو استماع، لأن تطابق الإجابات دليل قاطع على الغش، ويحتفظ بالأسئلة نفسها مستقبلًا مع تشديد العقوبة المعلنة، حتى يكون عبرة لغيرهم في السنوات القادمة.
  \item[\textbf{د)}] يلغي الاختبارات الإلكترونية نهائيًا ويعود إلى الاختبارات الورقية في كل المواد، لأن الاختبار الإلكتروني لا يمكن تصميمه بطريقة تقيس كفاءة الطالب بأي حال.
  \item[\textbf{هـ)}] يعالج الحالة وفق نظام المؤسسة بعد التحقق والاستماع للطلبة، ويعيد تصميم التقويم مستقبلًا بأسئلة تطبيق وتحليل مرتبطة بسياق الدرس، مع قواعد واضحة لما يُسمح به، أو بجزء حضوري شفهي.
\end{itemize}
\par\vspace{6pt}
\noindent
{\large\khatt\textbf{الوضعية 4}}\hfill{\small(الفصل 5)}\par\nopagebreak\vspace{3pt}
\setlength{\fboxsep}{7pt}\colorbox{gray!10}{\parbox{\dimexpr\linewidth-14pt}{طلب صحفي متدرّب من أداة استخراج أهم ثلاثة اقتباسات من حوار مسجّل نصه في عشرين صفحة مع روائي جزائري. جاءت الاقتباسات معبّرة، لكن أحدها لم يرد بحرفه في النص، بل جمع جملتين متباعدتين وأضاف إليهما ربطًا. ورئيس التحرير يريد نشر الحوار غدًا في الصفحة الأولى. ما التصرف المهني الأنسب؟}}\par\nopagebreak
\begin{itemize}[leftmargin=1.0cm, labelsep=6pt, itemsep=3pt, topsep=4pt, parsep=0pt]
  \item[\textbf{أ)}] يحذف علامتي التنصيص من الاقتباس المركّب ويُبقي نسبته الحرفية إلى الروائي في سياق الجملة، لأن غياب علامات التنصيص يكفي لجعل الدمج والإضافة مقبولين مهنيًا، ويحافظ بذلك على قوة الاقتباس وتأثيره في القراء.
  \item[\textbf{ب)}] يعتمد الاقتباسين الآخرين دون تحقق، لأن خطأ اقتباس واحد لا يعني احتمال أخطاء في غيره، والأداة تستخرج الجمل من النص المرفوع حرفيًا عادة، فيوفر وقته لمراجعة أجزاء أخرى من الحوار.
  \item[\textbf{ج)}] يعدّ الاقتباس المركّب غير صالح، لأن ما بين علامتي التنصيص يجب أن يطابق كلام الروائي حرفيًا؛ فيتحقق من كل اقتباس في النص حرفًا وسياقًا، ويستعمل الصياغة غير المباشرة إن أراد الجمع.
  \item[\textbf{د)}] يعتمد الاقتباس المركّب، لأنه يعبّر عن فكرة الروائي بدقة أكبر من الجملتين المتباعدتين، والدمج تحسين أسلوبي مقبول في الحوارات الصحفية المنشورة.
  \item[\textbf{هـ)}] يتخلى عن الأداة، ويكتب اقتباسات من ذاكرته بعد قراءة سريعة للحوار، لأن ذاكرة الصحفي أدق من أي أداة في نقل كلام المتحدث حرفيًا.
\end{itemize}
\par\vspace{6pt}
\noindent
{\large\khatt\textbf{الوضعية 5}}\hfill{\small(الفصل 6)}\par\nopagebreak\vspace{3pt}
\setlength{\fboxsep}{7pt}\colorbox{gray!10}{\parbox{\dimexpr\linewidth-14pt}{جرّبت أستاذة رفع صور إنشاءات تلاميذها المكتوبة بخط اليد إلى أداة متعددة الأنماط لاقتراح ملاحظات لغوية. كانت الملاحظات دقيقة مع الخطوط الواضحة، لكنها نسبت إلى أصحاب الخطوط الصعبة أخطاء إملائية لم يرتكبوها، لأنها أساءت قراءة كلماتهم. وتريد تعميم التجربة على أقسامها الأربعة إذا نجحت. ما التصرف الأنسب؟}}\par\nopagebreak
\begin{itemize}[leftmargin=1.0cm, labelsep=6pt, itemsep=3pt, topsep=4pt, parsep=0pt]
  \item[\textbf{أ)}] تتخلى عن الأداة كليًا حتى مع الخطوط الواضحة، لأن خطأها في بعض الأوراق يعني أن كل ملاحظاتها خاطئة، حتى الملاحظات الدقيقة التي تحققت منها بنفسها، وتعود إلى التصحيح اليدوي الكامل لكل الأوراق.
  \item[\textbf{ب)}] تتنبّه إلى أن أخطاء القراءة تُحمَّل ظلمًا لأصحاب الخطوط الصعبة، فتقابل ما قرأته الأداة بالورقة قبل اعتماد أي ملاحظة، أو تستعملها مع النصوص الواضحة فقط، لأن عدالة التقويم مسؤوليتها.
  \item[\textbf{ج)}] تخفّض علامات أصحاب الخطوط الصعبة بنسبة ثابتة لتعويض أخطاء الأداة، دون مقابلة ما قرأته بالورقة، لأن التعديل الجماعي أسرع وأعدل من المراجعة الفردية.
  \item[\textbf{د)}] تطلب من أصحاب الخطوط الصعبة إعادة كتابة إنشاءاتهم بالحاسوب في البيت، ثم تعتمد ملاحظات الأداة دون مراجعة، لأن الكتابة بالحاسوب تلغي كل أخطاء القراءة، فيتعودون على الكتابة الرقمية في الوقت نفسه.
  \item[\textbf{هـ)}] تعتمد الملاحظات كما هي، لأن الخط الصعب نفسه خطأ يستحق العقاب، والأخطاء الإملائية التي نسبتها الأداة إليهم جزاء مناسب على رداءة خطهم.
\end{itemize}
\par\vspace{6pt}
\begin{center}
{\khatt\textbf{جدول الإجابات}}\\[3pt]
\renewcommand{\arraystretch}{1.6}
\begin{tabular}{|>{\centering\arraybackslash}m{1.9cm}|*{5}{>{\centering\arraybackslash}m{1.2cm}|}}
\hline
\textbf{رقم السؤال} & \textbf{1} & \textbf{2} & \textbf{3} & \textbf{4} & \textbf{5} \\ \hline
\textbf{الإجابة} &  &  &  &  &  \\[8pt] \hline
\end{tabular}

\vspace{8pt}
\renewcommand{\arraystretch}{2.2}
\begin{tabular}{|>{\centering\arraybackslash}m{7.5cm}|>{\centering\arraybackslash}m{3.2cm}|>{\centering\arraybackslash}m{3.2cm}|}
\hline
\textbf{الاسم واللقب} & \textbf{الفوج} & \textbf{العلامة} \\ \hline
 & & \hfill /20 \\ \hline
\end{tabular}
\end{center}
\end{document}
